\section{Lattice ensemble and dispersion relations}
\label{sec:simulate}
We study the new momentum smearing method on 200 effectively decorrelated
configurations of $N_f=2$ non-perturbatively improved
Wilson fermions with the Wilson gluon action, generated by QCDSF
and RQCD. This constitutes a subset of
ensemble IV of Ref.~\cite{Bali:2014gha}.  Note that in
hadron structure studies we typically employ several sources
on about 2000 configurations~\cite{Bali:2014gha}, i.e.\
the present statistics are very moderate. Nevertheless, we
will report meaningful signals for momenta as high as
$2.8\units{GeV}$. The parameter values
of this $32^3\times 64$ ensemble read
$\beta=5.29$, $\kappa=0.13632$, corresponding to the inverse
lattice spacing  $a^{-1}\approx 2.76\units{GeV}$ and the (finite volume)
pion mass $m_{\pi}\approx 295\units{MeV}$. This gives a spatial
extent $L\approx 2.29\units{fm}\approx 3.42/m_{\pi}$.
Note that realizing the physical pion mass at this lattice spacing,
while keeping $L\gtrsim 3.5/m_{\pi}$, requires $L>70a$. As we are aiming
at momenta exceeding the hadron masses in any case, we do not
expect qualitative changes of results towards smaller pion masses
and have chosen the present parameters as a sufficiently
realistic but still affordable compromise. 

The momentum on a finite cubic lattice with even numbers of
points $N=L/a$ in each spatial direction
can take the discrete values
\begin{equation}
\label{quantize}
\vect{p}=\frac{2\pi}{L}\vect{P}\,,\quad
P_i\in\left\{-\frac{L}{2a}+1,-\frac{L}{2a}+2,\cdots, \frac{L}{2a}\right\}\,.
\end{equation}
This means the smallest non-trivial $|\vect{P}|=|(1,0,0)|$ gives
a momentum $|\vect{p}|\approx 0.54\units{GeV}$ while
$\vect{P}=(3,3,3)$ corresponds to $|\vect{p}|\approx 2.82\units{GeV}$.

The pion and nucleon masses have already been determined with high
statistics in Ref.~\cite{Bali:2014nma} and read
\begin{equation}
\label{eq:masses}
am_{\pi}=0.10675(52)\,,\quad am_N=0.3855(45)\,.
\end{equation}
We will compare our pion and nucleon
ground state energies to expectations from continuum and lattice
dispersion relations, using these reference values.
The continuum dispersion relation reads
\begin{equation}
\label{eq:disperse}
aE=a\sqrt{m^2+\vect{p}^2}\,.
\end{equation}
In addition, we will compare the pion energies to
the lattice dispersion relation for
a free naively discretized scalar particle,
\begin{equation}
\label{eq:dispersepi}
\cosh(aE_{\pi})=\cosh(am_{\pi})+\frac{a^2}{2}\hat{\vect{p}}^2\,,\quad
\end{equation}
and the nucleon energies to the dependence expected 
for a free Wilson fermion with the Wilson parameter $r=1$, see,
e.g., Ref.~\cite{ElKhadra:1996mp},
\begin{equation}
\label{eq:disperseN}
\cosh(aE_N)=1+\frac{\left(e^{am_N}-1+a^2\hat{\vect{p}}^2/2\right)^2+a^2\overline{\vect{p}}^2}{2\left(e^{am_N}+a^2\hat{\vect{p}}^2/2\right)}\,.
\end{equation}
Above, we used the standard abbreviations
\begin{equation}
\hat{p}_{\mu}=\frac{2}{a}\sin\left(\frac{ap_{\mu}}{2}\right)\,,\quad
\overline{p}_{\mu}=\frac{1}{a}\sin(ap_{\mu})\,,
\end{equation}
and $\hat{\vect{p}}^2=\sum_j\hat{p}_j^2$,
$\overline{\vect{p}}^2=\sum_j\overline{p}_j^2$.
We remark that the mass parameters $m_0$ within the
naive propagators (e.g., for the scalar case:
$1/(m_0^2+\hat{p}_{\mu}\hat{p}_{\mu})$)
differ from the rest frame energies by lattice artefacts.
The above masses are defined to satisfy $m=E(\vect{0})$, as it should be.
Their conversions to
$m_0$ are given as
$m_{\pi}=2a^{-1}\sinh (am_0/2)$ and $m_N=a^{-1}\ln(1+am_0)$, respectively.

\begin{figure}[ht]
\centerline{\includegraphics[width=0.95\textwidth]{smear3d.pdf}}
\caption{Cross sections of the smearing density profile $\rho(\vect{x})$
Eq.~\protect\eqref{density} in the $x_1$-$x_2$ plane.
The colour encodes the phase Eq.~\protect\eqref{phase}
and the momentum smearing $\vect{k}$ parameter has the direction $(1,1,0)$.
Top left: free field case. Top right: APE smeared gauge links.
Bottom left: original gauge links. Bottom right: APE smeared
links with an additional boost factor $\gamma=5.3$.}
\label{fig:smearing}
\end{figure}

Obviously, the continuum dispersion relation should become
violated towards large momenta. In this case, we would not expect the
lattice dispersion relations, that apply to free pointlike particles,
to accurately describe the data either. However, the difference between
the continuum and the lattice formulae can serve as a naive estimate
of the expected size of lattice artefacts.

